\subsection{初等滤子}

定义 7. 设 \((x_{n})_{n\in \mathbb{N}}\) 是集合 \(\mathbf{X}\) 的元素的无穷序列。与序列 \((x_{n})\) 相伴的初等滤子是由 Fréchet 滤子（第 1 目）经映射 \(n\to x_n\)（\(\mathbf{N}\) 到 \(\mathbf{X}\) 中的）所得的像生成的滤子。等价地说，与 \((x_{n})\) 相伴的初等滤子是子集 \(\mathbf{M}\)（\(\mathbf{X}\) 的，使得都有 \(x_{n}\in \mathbf{M}\)——除去 \(n\) 的有限多个值外）所成之集。若用 \(\mathbf{S}_n\) 表示所有这样的 \(x_{p}\)（满足 \(p\geqslant n\)）所成之集，则诸集合 \(\mathbf{S}_n\) 组成与序列 \((x_{n})\) 相伴的初等滤子的一个基。

与序列 \((x_{n})\) 的无穷子列相伴的初等滤子细于与 \((x_{n})\) 相伴的初等滤子（参看习题 15）。

按定义，每个初等滤子都有可数基。反之：命题 11. 若滤子 F 有可数基，则它是细于 F 的诸初等滤子之交。

把 \(\mathfrak{F}\) 的可数基排成序列 \((\mathbf{A}_n)_{n\in \mathbb{N}}\)；若令

\[
\mathrm{B} _ {n} = \bigcap_ {p = 0} ^ {n} \mathrm{A} _ {p},
\]

则诸 \(B_{n}\) 仍组成 F 的基（第 3 目，命题 3），且对每个 n 有 \(B_{n+1} \subset B_{n}\)。设 \(a_{n}\) 是 \(B_{n}\) 的任一元素（对每个 \(n \in N\)）；则显然 F 粗于与 \((a_{n})\) 相伴的初等滤子。于是细于 F 的诸初等滤子之交 S 存在且细于 F；若 S 严格细于 F，则存在集合 \(M \in S\)，使对每个 n 有 \(B_{n} \cap CM \neq \varnothing\)；若 \(b_{n} \in B_{n} \cap CM\)，则与序列 \((b_{n})\) 相伴的初等滤子细于 F 且不含 M。这与 S 的定义矛盾。

注. 粗于某个具有可数基的滤子的滤子未必具有可数基；例如，若 X 是不可数无限集，则由 X 的有限子集的余集组成的滤子没有可数基（否则 X 的有限子集所成之集将是可数的，与假设相反）；然而这个滤子粗于由 X 的互异元素组成的无穷序列所相伴的每个初等滤子。

